% -------------------------------------------------------------------------------------
\section{SCET\texorpdfstring{$_{\rm II}$}{II}、零箱减除与快度发散}\label{sec:scet-scet2}

有效理论中的共线圈积分在圈动量趋于软（或超软）区域时与软圈积分重叠；这正是第~9~章减除法要处理的重复计算。

\begin{definition}[零箱减除]\label{def:scet-zero-bin}
(a) 设 $I_c(k)$ 是 SCET 共线图的被积函数（$k$ 为共线圈动量或实辐射动量）。其\emph{零箱}（zero-bin）$I_c^{(0)}(k)$ 是把 $k$ 按嵌套在其中的较小区域（软或超软模式）重新计数后 $I_c$ 的领头项；物理的共线贡献定义为 $\tilde I_c=\int(I_c-I_c^{(0)})$，多个嵌套区域时按包含关系逐级减除\cite{ManoharStewart2007}。这是\cref{def:method-subtraction}中 $C_R\Gamma=T_R(\Gamma-\sum_{R'<R}C_{R'}\Gamma)$ 在有效理论中的对应：$T_R$ 是共线近似（SCET 的共线 Feynman 规则），$R'$ 是嵌套于 $R$ 的软区域。
(b) 当红外发散也用维数正规化时，零箱减除与“除以软 Wilson 线矩阵元”的软减除等价；这一点在夸克形状因子与 $x_B\to1$ 的 DIS 中、并在两圈阿贝尔虚修正中被显式检验\cite{IdilbiMehen2007}；两种形式中消除重复计数的比较见\cite{Lee:2006nr}。只有当零箱减除在每个共线因子中自洽地实施时，重复计数才被避免。
(c) 对推力的单圈喷注函数（\cref{sec:scet-rg}），实辐射零箱的被积函数为 $\propto\delta^+(k^2)\,\frac{1}{n\cdot k\,\bar n\cdot k}\,\delta\big(n\cdot k-p^2/\bar n\cdot p\big)$（测量只约束 $n\cdot k$）。在 $\bar n\cdot k\to a\,\bar n\cdot k$、$k_\perp\to\sqrt a\,k_\perp$、$n\cdot k$ 固定之下，$\delta(k^2)$ 给出 $a^{-1}$，$1/\bar n\cdot k$ 给出 $a^{-1}$，测度 $\dd(\bar n\cdot k)\,\dd^{2-2\varepsilon}k_\perp$ 给出 $a^{2-\varepsilon}$，故零箱积分在 $a$ 下是 $-\varepsilon$ 次齐次的，从而无标度，按\cref{lem:prelim-scaleless}为零；虚修正的零箱同样无标度。于是推力的减除后喷注函数等于未减除的喷注函数。% w18_s4 (D); App N prop:appN-thrust-zerobin is the Collins-form counterpart (pointer)
\uses{def:method-subtraction, lem:prelim-scaleless}
\end{definition}

在 Collins 减除的语言中，(c) 的对应陈述是\cref{prop:appN-thrust-zerobin}（只依赖于发散链项的取值约定）。

\paragraph{SCET$_{\rm II}$。}
在横动量分布、$B$ 介子的遍举衰变等问题中，与共线模式 $(1,\lambda^2,\lambda)$ 发生相互作用的是软模式 $(\lambda,\lambda,\lambda)$，二者虚度相同。

\begin{proposition}[SCET$_{\rm II}$ 与软 Wilson 线 \tagR]\label{prop:scet-scet2}
（$R_{\mathrm{PT}}$：在重整化微扰论中对红外正规化的部分子态严格。）（树图层次。）(a) 共线模式与软模式的虚度都是 $\lambda^2Q^2$；软动量 $k$ 被 $n$ 共线线吸收（或由其发射）后，中间线的虚度为 $(p+k)^2\simeq\bar n\cdot p\,n\cdot k\sim\lambda Q^2$，是硬共线的。(b) 在树图层次积掉这些硬共线线后，软胶子只通过算符中的软 Wilson 线
\begin{equation}\label{eq:scet-Sn}
S_n(x)=\mathrm P\exp\Big(\ii g\int_{-\infty}^0\dd s\;n\cdot A_s(x+sn)\Big)
\end{equation}
与 $n$ 共线夸克耦合（例如两喷注流化为 $\bar\chi_nS_n^\dagger\,\Gamma\,S_{\bar n}\chi_{\bar n}$），SCET$_{\rm II}$ 的领头阶拉氏量中没有软—共线耦合。(c) 软贡献与共线贡献只在快度上分开，单个 SCET$_{\rm II}$ 积分可以有维数正规化不能正规化的\emph{快度发散}，必须引入快度正规子\cite{Hill:2002vw,BauerEtAl2002}。（本命题只涉及部分子振幅。）
\uses{def:scet-modes, prop:scet-wilson}
\end{proposition}

\begin{proof}
(a) 由 $2p\cdot k=\bar n\cdot p\,n\cdot k+n\cdot p\,\bar n\cdot k+2p_\perp\cdot k_\perp$，三项分别为 $\lambda$、$\lambda^3$、$\lambda^2$（以 $Q^2$ 为单位），而 $k^2\sim\lambda^2$。(b) 这与\cref{prop:scet-wilson}证明的第二、三步相同，只是 $n$ 与 $\bar n$ 互换：对出射的 $n$ 共线夸克，每一对“硬共线传播子 × 顶角”在领头幂次化为 $\bar u(p)\,\ii g\,\sla\epsilon^*t^a\,\frac{\ii\frac{\bar n\cdot p}{2}\sla n}{\bar n\cdot p\,(n\cdot k+\ii0)}=-g\,t^a\frac{n\cdot\epsilon^*}{n\cdot k+\ii0}\,\bar u(p)$（用 $\bar u_n\sla n=0$），$r$ 个胶子按次序求和给出 $S_n^\dagger$ 的展开（对 $n\cdot k\neq0$，$\ii0$ 的取法无关，与\cref{rem:scet-y-direction}相同）。对一个软胶子，精确的发射振幅与程函因子 $-n\cdot\epsilon^*/n\cdot k$ 之差（对夸克自旋求和后）的相对大小在一个代表性运动学点上为 $0.81$、$0.081$、$0.0081$（$\lambda=10^{-1},10^{-2},10^{-3}$），按 $\lambda$ 的一次方减小；具体数值依赖于所选的点，标度行为不依赖。% w18_s5 (H); c_scet quotes 0.78, 0.077, 0.0077 for a single spinor at the same point
软胶子若出现在拉氏量中与共线场耦合，就会把共线粒子推到硬共线虚度，因此在 SCET$_{\rm II}$ 中它们只能出现在算符的 $S_n$ 里。(c) 取一个软实辐射，其横动量 $\bm k_{\mathrm T}$ 被测量固定。令 $n\cdot k=k_T\ee^{-y}$、$\bar n\cdot k=k_T\ee^{y}$，
\begin{equation}\label{eq:scet-rapdiv}
\int\dd(n\cdot k)\,\dd(\bar n\cdot k)\,\frac{\delta(n\cdot k\,\bar n\cdot k-\bm k_{\mathrm T}^2)}{n\cdot k\,\bar n\cdot k}=\frac{1}{\bm k_{\mathrm T}^2}\int_{-\infty}^{\infty}\dd y ,
\end{equation}
在 $y\to\pm\infty$ 处发散；由于横向积分已被测量固定，$D$ 维测度对它不起作用。
\end{proof}

快度正规子有三类：Collins 的非类光 Wilson 线（第~14~章，\cref{def:tmd-collins}）；Chiu--Jain--Neill--Rothstein（CJNR）的 $\eta$ 正规子与快度重整化群\cite{ChiuEtAl2012a,ChiuEtAl2012b}，它在软 Wilson 线的每个顶点插入 $w\,|2k^z/\nu|^{-\eta/2}$（于是单胶子实辐射带有 $w^2\nu^\eta|2k^z|^{-\eta}$），在共线 Wilson 线的每个顶点插入 $w\,|\bar n\cdot k/\nu|^{-\eta}$；Becher--Neubert 的解析正规子与“共线反常”\cite{BecherNeubert2011}。以 $\eta$ 正规子为例，$|2k^z|=2k_T|\sinh y|$，\cref{eq:scet-rapdiv}中的快度积分变为
$\int_{-\infty}^{\infty}\dd y\,|2\sinh y|^{-\eta}=2^{-\eta}\Gamma(\tfrac\eta2)\Gamma(\tfrac12-\tfrac\eta2)/\sqrt\pi=\tfrac2\eta+\order{\eta}$：一个 $1/\eta$ 极点同时正规化了两端；软辐射在正、负快度两侧都须截断，这是 CJNR 在软 Wilson 线中取 $\eta/2$ 的物理原因\cite{ChiuEtAl2012a}。% w18_s5 (A2, A3)

\begin{proposition}[快度发散的处理与等价性 \tagA]\label{prop:scet-rapidity-rg}
对横动量依赖（TMD）的因子化：
(a) 快度重整化群（CJNR）、共线反常（Becher--Neubert）与 Collins 的非类光 Wilson 线（\cref{def:tmd-collins}）给出等价的 TMD 因子化\cite{Collins:2017oxh}：在每一种方案中，单个共线（束流）函数与软函数依赖于正规子，而进入截面的组合与之无关。
(b) 支配 $b$ 空间截面在固定 $b_T$ 下 $Q^2$ 依赖的演化核在三种方案中的对应为
\begin{equation}\label{eq:scet-kernels}
\tilde K(b_T;\mu)=\tfrac12\,\tilde\gamma_\nu^S(b_T;\mu)=-F_{q\bar q}(b_T^2,\mu)=-2D(b_T;\mu),
\end{equation}
其中 $\tilde K$ 是\cref{prop:tmd-evolution}的 Collins--Soper 核；$\tilde\gamma_\nu^S$ 由 $\nu\frac{\dd}{\dd\nu}\tilde S=\tilde\gamma^S_\nu\tilde S$ 定义（CJNR 的 $\eta$ 正规子，夸克 Wilson 线，且 $\tilde\gamma^S_\nu+2\tilde\gamma^B_\nu=0$）；$F_{q\bar q}$ 由 $[\tilde B_{q}\tilde B_{\bar q}]_{Q^2}=(b_T^2Q^2/b_0^2)^{-F_{q\bar q}}B_qB_{\bar q}$ 定义（Becher--Neubert）；$D$ 是 $\delta$ 正规子（Echevarria--Idilbi--Scimemi）归一化中 $\partial\ln\tilde f/\partial\ln Q^2=-D$（$\zeta=Q^2$）的指数\cite{EchevarriaIdilbiScimemi2012}。单圈时四者都等于 $-\frac{\as C_F}{\pi}L_b$，$L_b=\ln(\mu^2b_T^2/b_0^2)$。
\uses{prop:scet-scet2, def:tmd-collins, prop:tmd-evolution}
\end{proposition}

\begin{proofsketch}
(1) 在每种方案中，束流函数与软函数的快度发散在二者的乘积中相消（CJNR：$\tilde\gamma^S_\nu+2\tilde\gamma^B_\nu=0$\cite{ChiuEtAl2012b}；Becher--Neubert：$Q^2$ 依赖被完整地指数化\cite{BecherNeubert2011}；Collins：\cref{def:tmd-collins}中的平方根组合）。(2) 固定 $b_T$ 时，$b$ 空间截面 $H\,\tilde f_A\tilde f_B$ 对 $\ln Q^2$ 的导数是物理量。在 Collins 方案中取 $\zeta_A=\zeta_B=Q^2$，它等于 $\partial\ln H/\partial\ln Q^2+\tilde K$（\cref{prop:tmd-evolution}）；在 $\eta$ 方案中 $\tilde f=\tilde B\sqrt{\tilde S}$，$\tilde B$ 只通过 $\nu/\omega$ 依赖于大光锥动量 $\omega$，故 $\partial\ln\tilde f/\partial\ln\omega=-\tilde\gamma^B_\nu=\frac12\tilde\gamma^S_\nu$，而 $\zeta\propto\omega^2$；在 Becher--Neubert 与 $\delta$ 方案中，乘积分别正比于 $(Q^2)^{-F_{q\bar q}}$ 与 $(Q^2)^{-2D}$。硬函数相同时，比较即得\cref{eq:scet-kernels}；RG 一致性又给出 $\frac{\dd\tilde K}{\dd\ln\mu}=-2\Gamma_{\rm cusp}$、$\frac{\dd F_{q\bar q}}{\dd\ln\mu}=2\Gamma_{\rm cusp}$、$\frac{\dd D}{\dd\ln\mu}=\Gamma_{\rm cusp}$（$\Gamma_0=4C_F$）。(3) 单圈显式核对：$\tilde K=-\frac{\as C_F}\pi L_b$（第~14~章）；CJNR 的 $p_T$ 空间结果 $\nu\frac{\dd}{\dd\nu}S^{(1)}=\frac{4\as C}{\pi}\mathcal L_0$\cite{ChiuEtAl2012b}（$C$ 为 Wilson 线的 Casimir，$\mathcal L_0=\frac1{2\pi}[1/\bm p^2]_+$ 在 $|\bm p|<\mu$ 内减除），用 $\int\dd^2p\,\ee^{\ii\bm p\cdot\bm b}\mathcal L_0=\int_0^\infty\frac{\dd p}p[J_0(pb)-\theta(\mu-p)]=-\frac12L_b$ 变到 $b$ 空间，得 $\tilde\gamma^S_\nu=-\frac{2\as C_F}{\pi}L_b$；Becher--Neubert 的 $F_{q\bar q}=\frac{C_F\as}\pi L_\perp$（$L_\perp=L_b$）\cite{BecherNeubert2011}；Echevarria--Idilbi--Scimemi 的 $D=\frac{\as}{4\pi}\,2C_FL_T$（$L_T=L_b$）\cite{EchevarriaIdilbiScimemi2012}。四者满足\cref{eq:scet-kernels}。% w18_s5 (B, C): Fourier transform and the one-loop dictionary checked
完整论证见\cite{Collins:2017oxh}（各种表述之间的联系）与\cite{ChiuEtAl2012b,BecherNeubert2011,EchevarriaIdilbiScimemi2012}。
\end{proofsketch}

\paragraph{地位与缺口}
\cref{prop:scet-rapidity-rg}是全阶论证（[A]）：各种方案的等价性是逐阶、并在低阶显式验证的（据文献\cite{Collins:2017oxh}摘要，它推导了 CSS 原始表述、Collins 表述与若干 SCET 表述之间的联系，并以 $\order{\as^2}$、$\order{\as^3}$ 的已知结果核对），没有一个对所有阶、对非微扰函数同时成立的证明；其中每一种方案本身都以 TMD 因子化定理（\cref{thm:tmd-factorization}，[A]）为前提，后者的瓶颈是\cref{asm:appN-tmd}与 Glauber 区（\cref{asm:appO-mixing}）。最弱的一步是 (a) 的“等价”本身：单个束流函数与软函数依赖于正规子，只有组合是物理的——快度发散标志着一个真实的附加因子化变量（$\zeta$ 或 $\nu$），而不是某种正规子的人为产物；这一步是方案之间的等价性问题，不涉及幂次误差。TMD 算符的多重可重整性在文献中被论证而未被证明\cite[第~13~章]{Collins2011}。升级途径：在一个固定的方案中把多重可重整性（\cref{asm:appN-tmd}(T3)）证明到所有阶，并给出方案之间的有限重定义。
